\chapter{Disjoint Sets and Union-Find}

In the previous chapter, B-trees organized massive ordered dictionaries across secondary storage pages, where queries drilled down an ordered key space. But many computational problems care nothing about key order. You are given a pile of objects -- people, computers on a network,
pixels in a photograph, cells in a maze -- and from time to time someone tells
you that two of them belong together. ``Alice and Bob are friends.'' ``Cable
runs from router 3 to router 7.'' ``This pixel touches that pixel and they are
the same color.'' As these facts arrive, the objects clump up into groups. And
at any moment, someone might stop you and ask: \emph{are these two particular
objects currently in the same group?}

That is the entire problem. You maintain a collection of groups, you
occasionally merge two groups into one, and you occasionally ask whether two
objects share a group. It is called the \emph{disjoint-set} or
\emph{union-find} problem, and its best solution is one of the most surprising
results in data structures: an implementation so fast that bounding its
per-operation cost requires the extraordinarily slow-growing inverse Ackermann function.

\begin{intuitionbox}
Picture a room full of strangers. Each begins alone, a group of one. You
walk around announcing introductions: ``you two know each other now.'' When
two people meet, their entire friend-groups silently merge -- if Alice knew
Carol and Bob knew Dave, then after Alice meets Bob, all four are one group.
At any point I can tap two people on the shoulder and ask, ``same group?'' To
answer instantly, each group elects a single \emph{leader}, and everyone in
the group knows (perhaps after a short chain of pointers) who their leader is.
Two people are in the same group exactly when they share a leader. That single
idea -- \emph{a group is named by its leader} -- is the whole chapter.
\end{intuitionbox}

We will start with the slowest reasonable version, watch exactly where it
hurts, and add two small tricks. Those two tricks together turn a structure
that could cost $\Theta(n)$ per operation into one that costs, for all
practical purposes, a constant. Then we will meet the strange, slow-growing
function that measures precisely how close to constant we got.

\section{What We Are Actually Building}
\label{ch19:problem}

Let us pin the problem down. We have a universe of $n$ elements. They are
partitioned into \emph{disjoint sets} -- ``disjoint'' meaning no element lives
in two sets at once, so every element belongs to exactly one set. We want to
support three operations.

\begin{formalbox}
A \textbf{disjoint-set data structure} maintains a collection
$\mathcal{S} = \{S_1, S_2, \dots, S_k\}$ of disjoint dynamic sets. Each set is
identified by a \emph{representative}, which is some fixed member of the set.
Three operations are supported:
\begin{itemize}
  \item $\textsc{Make-Set}(x)$: create a new set whose only member (and
        representative) is $x$. Since $x$ joins no other set, the sets stay
        disjoint.
  \item $\textsc{Union}(x, y)$: replace the two sets containing $x$ and $y$,
        say $S_x$ and $S_y$, by their union $S_x \cup S_y$, and destroy the two
        originals. The new set gets one representative (often a member of one
        of the two old sets).
  \item $\textsc{Find-Set}(x)$: return the representative of the set
        containing $x$.
\end{itemize}
\end{formalbox}

Notice how cleanly $\textsc{Find-Set}$ answers our shoulder-tap question.
``Are $x$ and $y$ in the same set?'' is just the test
$\textsc{Find-Set}(x) = \textsc{Find-Set}(y)$. We never need a separate
``same-set'' operation; representatives give it to us for free. The only
promise the representative must keep is \emph{stability}: if you call
$\textsc{Find-Set}$ twice on the same element with no intervening
$\textsc{Union}$, you must get the same answer both times. Which element
serves as representative is up to us, and we will exploit that freedom
ruthlessly.

\begin{focusbox}
Hold these three operations in your head as the contract. Everything that
follows -- linked lists, forests, ranks, path compression -- is just
\emph{implementation} behind that contract. The user only ever sees
$\textsc{Make-Set}$, $\textsc{Union}$, and $\textsc{Find-Set}$.
\end{focusbox}

\subsection{A motivating use: connected components}
\label{ch19:components}

Before we build anything, here is why the contract is worth keeping. A
\emph{graph} is a set of vertices and a set of edges joining them; two vertices
are in the same \emph{connected component} if you can walk from one to the
other along edges. Suppose edges arrive one at a time and, after each, you want
to know how many components remain, or whether two given vertices are connected
yet. Union-find answers this beautifully.

\begin{algorithm}[H]
\caption{\textsc{Connected-Components}$(G)$}
\KwIn{A graph $G = (V, E)$.}
\KwOut{A disjoint-set structure where $u, v$ share a set iff they are connected.}
\ForEach{vertex $v \in V$}{
  $\textsc{Make-Set}(v)$\;
}
\ForEach{edge $(u, v) \in E$}{
  \If{$\textsc{Find-Set}(u) \neq \textsc{Find-Set}(v)$}{
    $\textsc{Union}(u, v)$\;
  }
}
\Return the disjoint-set structure\;
\end{algorithm}

After this runs, the query procedure \textsc{Same-Component}$(u,v)$ is just one line:
return whether $\textsc{Find-Set}(u) = \textsc{Find-Set}(v)$. We will meet this
exact pattern again in Chapter~23, where Kruskal's minimum-spanning-tree
algorithm leans entirely on union-find to avoid creating cycles. So this is not
a toy: it is an indispensable pattern you will reuse.

\begin{aestheticbox}
Union-find quietly powers a surprising amount of real software. Kruskal's
minimum-spanning-tree algorithm uses it to detect whether adding an edge would
form a cycle. Image-processing tools use it for \emph{connected-component
labeling} -- finding the blobs in a black-and-white image. Compilers use it for
\emph{unification} in type inference, merging type variables that must be
equal. Maze generators carve passages by unioning cells. Even physics
simulations of percolation (does water seep through a porous rock?) are, at
heart, union-find. It is the data structure you did not know you were already
using.
\end{aestheticbox}

\section{First Try: a Linked-List Representation}
\label{ch19:linkedlist}

Let us build the simplest thing that could possibly work, so we can see exactly
where its weaknesses bite. Represent each set as a \emph{linked list}. The first
object in the list is the representative. Each object stores a pointer to the
next object in its list, and -- this is the clever part -- also a pointer
straight back to the representative.

\begin{center}
\begin{tikzpicture}[
  node distance=12mm,
  obj/.style={draw, rounded corners, minimum width=8mm, minimum height=7mm,
    fill=blue!4},
  rep/.style={obj, fill=green!10, very thick},
  >={Stealth[length=2.2mm]}]
  \node[rep] (a) {$c$};
  \node[obj, right=of a] (b) {$h$};
  \node[obj, right=of b] (c) {$e$};
  \node[obj, right=of c] (d) {$b$};
  \draw[->] (a) -- (b);
  \draw[->] (b) -- (c);
  \draw[->] (c) -- (d);
  \draw[->, dashed, blue!55!black, bend right=38] (b) to (a);
  \draw[->, dashed, blue!55!black, bend right=45] (c) to (a);
  \draw[->, dashed, blue!55!black, bend right=50] (d) to (a);
  \node[below=10mm of b, font=\small, gray] {dashed = ``back to representative'' pointer};
\end{tikzpicture}
\end{center}

The solid arrows chain the list together (so we can walk it); the dashed arrows
let every member point home to the representative $c$. With this layout,
$\textsc{Make-Set}$ and $\textsc{Find-Set}$ are both exceptionally fast.

\begin{itemize}
  \item $\textsc{Make-Set}(x)$: build a one-element list. Constant time.
  \item $\textsc{Find-Set}(x)$: follow $x$'s back-pointer to the
        representative. Constant time.
\end{itemize}

The trouble is $\textsc{Union}$. To union two lists, we append one onto the
other. Appending the list is cheap -- splice one tail onto the other. But every
object on the appended list now has the \emph{wrong} back-pointer: it still
points to its old representative. We must walk the appended list and repoint
each one to the new representative. That walk costs time proportional to the
length of the list we appended.

\begin{trapbox}
The expensive part of the linked-list union is not joining the lists -- that is
one pointer splice. It is \emph{updating every back-pointer} on the list that
got absorbed. Beginners often assume union is cheap because ``it's just a list
append.'' The back-pointers are the hidden tax, and they are what make this
representation slow.
\end{trapbox}

\subsection{How bad can it get? A worst-case sequence}
\label{ch19:linkedlist-worst}

Naively, $\textsc{Union}(x,y)$ might always append $y$'s list onto $x$'s. An
adversary can exploit this. Suppose we make $n$ singleton sets
$x_1, \dots, x_n$, then perform $n-1$ unions, always appending the long
accumulated list onto a fresh singleton:

\[
\textsc{Union}(x_2, x_1),\ \textsc{Union}(x_3, x_2),\ \dots,\
\textsc{Union}(x_n, x_{n-1}).
\]

\begin{center}
\begin{tikzpicture}[
  bar/.style={draw, fill=trapred!20, minimum height=4mm},
  >={Stealth[length=2mm]}, font=\small]
  \foreach \k/\h in {1/1,2/2,3/3,4/4,5/5}{
    \node[bar, minimum width=\h*4mm, anchor=west] at (0, -\k*6mm) {};
    \node[anchor=east] at (-1mm, -\k*6mm) {union \k};
    \node[anchor=west] at (\h*4mm+2mm, -\k*6mm) {repoint \h};
  }
  \node[anchor=west] at (0, -6*6mm) {$\vdots$};
  \node[font=\footnotesize, gray, anchor=west] at (0,-7.4*6mm)
    {each union absorbs the growing list, so the work climbs $1,2,3,\dots,n-1$};
\end{tikzpicture}
\end{center}

Watch the cost if each union repoints the \emph{growing} side. The first union
repoints $1$ object, the second repoints $2$, the third $3$, and so on, so the
total update work across all $n-1$ unions is
\[
1 + 2 + 3 + \cdots + (n-1) = \frac{(n-1)\,n}{2} = \Theta(n^2).
\]
Averaged over the $2n-1$ operations in the sequence, that is $\Theta(n)$ per
operation. For a million elements, a single union near the end could touch a
million back-pointers. That is the pain point. Now let us fix it.

\subsection{The weighted-union heuristic}
\label{ch19:weighted}

The fix is embarrassingly simple, and it is the first appearance of a theme
that will repeat: \emph{when you merge two things, always absorb the smaller
into the larger}. Keep, with each list, a count of its length. When you union,
append the shorter list onto the longer one, and repoint only the shorter
list's back-pointers. This is the \emph{weighted-union heuristic}.

Why does this help so dramatically? Think about a single object $x$, and ask:
how many times can $x$'s back-pointer ever be updated? Its back-pointer is
updated only when $x$ is on the \emph{smaller} of two lists being merged. But
the moment that happens, $x$ lands in a merged list at least \emph{twice} as
big as the one it left. So each time $x$ gets repointed, the size of $x$'s set
at least doubles.

\begin{intuitionbox}
A quantity that at least doubles each time something happens cannot have that
something happen very often. Start at size $1$; after $k$ doublings the set has
size at least $2^k$. Since no set exceeds $n$, we need $2^k \le n$, so
$k \le \log_2 n$. Therefore \emph{any single object's back-pointer is updated
at most $\log_2 n$ times over the entire history of the structure.}
\end{intuitionbox}

Summing across all $n$ objects, the total back-pointer work over any sequence
of operations is at most $n \log_2 n$. The earlier quadratic blow-up is gone.

\begin{formalbox}
\textbf{Theorem (weighted union).} Using the linked-list representation with
the weighted-union heuristic, any sequence of $m$ \textsc{Make-Set},
\textsc{Union}, and \textsc{Find-Set} operations, of which $n$ are
\textsc{Make-Set}, takes $O(m + n \log n)$ time in total.
\end{formalbox}

This is a good result -- it is exactly the bound CLRS proves for this
representation. But we can do better still. A second way to picture sets opens
the door to near-constant time: we trade lists for trees.

\section{The Forest Representation}
\label{ch19:forest}

Here is the idea that carries us the rest of the way. Represent each set not as
a list but as a \emph{rooted tree}. Each element has a single pointer to its
\emph{parent}. The root of a tree points to itself (a convenient way to say ``I
am the representative''). The whole structure is therefore a \emph{forest} -- a
collection of these trees, one per set. The representative of a set is the root
of its tree.

\begin{center}
\begin{tikzpicture}[
  level distance=11mm, sibling distance=15mm,
  every node/.style={draw, circle, minimum size=7mm, fill=blue!4,
    inner sep=0pt, font=\small},
  >={Stealth[length=2mm]}, edge from parent/.style={draw,->}]
  \node[fill=green!12, very thick] {$c$}
    child {node {$h$}}
    child {node {$e$}
      child {node {$b$}}
    };
\end{tikzpicture}
\hspace{1.4cm}
\begin{tikzpicture}[
  level distance=11mm, sibling distance=12mm,
  every node/.style={draw, circle, minimum size=7mm, fill=blue!4,
    inner sep=0pt, font=\small},
  >={Stealth[length=2mm]}, edge from parent/.style={draw,->}]
  \node[fill=green!12, very thick] {$f$}
    child {node {$g$}}
    child {node {$d$}};
\end{tikzpicture}
\end{center}

Each arrow points from a child \emph{up} to its parent. These are not
binary trees and not search trees. There is no ordering, no balance condition
yet, and a node may have any number of children. The only thing that matters is
\emph{which root you reach by walking up}.

The three operations become:

\begin{itemize}
  \item $\textsc{Make-Set}(x)$: make a one-node tree. Set $x.p = x$ (its parent
        is itself). Constant time.
  \item $\textsc{Find-Set}(x)$: walk up parent pointers from $x$ until you hit a
        node that is its own parent. That node is the root; return it.
  \item $\textsc{Union}(x, y)$: find the two roots, then make one root point to
        the other. One pointer change (plus two finds).
\end{itemize}

Here is the bare-bones pseudocode, before any cleverness.

\begin{algorithm}[H]
\caption{Naive forest operations}
\SetKwProg{Fn}{function}{}{}
\Fn{\textsc{Make-Set}$(x)$}{
  $x.p \gets x$\;
}
\Fn{\textsc{Find-Set}$(x)$}{
  \While{$x \neq x.p$}{
    $x \gets x.p$\;
  }
  \Return $x$\;
}
\Fn{\textsc{Union}$(x, y)$}{
  $\textsc{Link}\big(\textsc{Find-Set}(x), \textsc{Find-Set}(y)\big)$\;
}
\Fn{\textsc{Link}$(u, v)$}{
  \tcp{$u$ and $v$ are roots}
  $u.p \gets v$\;
}
\end{algorithm}

The cost of everything is the cost of $\textsc{Find-Set}$, which is the cost of
walking from a node to its root -- that is, the \emph{height} of the tree. And
here is the catch: nothing yet stops the trees from growing tall and stringy.
If $\textsc{Link}$ always hangs the first root under the second, the same
adversarial union chain from before builds a single path of $n$ nodes, and
$\textsc{Find-Set}$ on the deepest node walks all $n$ of them. We are right back
to $\Theta(n)$.

\begin{walkbox}
Before reading on, take this on a walk. The linked-list version had a
weighted-union trick: absorb the smaller into the larger. The forest version is
stringy because we attach roots carelessly. What is the analogous ``attach the
smaller thing under the larger thing'' rule for trees -- and what should
``smaller'' even mean when the thing is a tree, not a list? There are two
natural answers, and they behave almost identically. See if you can invent at
least one of them.
\end{walkbox}

\section{Trick One: Union by Rank}
\label{ch19:rank}

The trees got tall because we attached roots blindly. The cure is the same
philosophy as weighted union: \emph{when you link two roots, put the shorter
tree under the taller one's root.} If you always hang the shorter tree beneath
the taller, the combined tree's height does not grow -- the shorter tree's
nodes are now one level deeper, but they were shallower to begin with, so they
still fit under the old height. The total height only increases when you link
two trees of \emph{equal} height, and even then it goes up by just one.

We could track exact heights, but maintaining exact heights becomes awkward
once we add the second trick (path compression shortens paths during queries
without recomputing subtree heights). So instead we track an \emph{upper bound} on the
height called the \emph{rank}. Every node carries an integer $x.\textit{rank}$.
A fresh singleton has rank $0$. When we link two roots:

\begin{itemize}
  \item if their ranks differ, hang the lower-rank root under the higher-rank
        root, and \emph{no rank changes};
  \item if their ranks are equal, pick either as the new root and increment its
        rank by one.
\end{itemize}

\begin{algorithm}[H]
\caption{Forest with union by rank}
\SetKwProg{Fn}{function}{}{}
\Fn{\textsc{Make-Set}$(x)$}{
  $x.p \gets x$\;
  $x.\textit{rank} \gets 0$\;
}
\Fn{\textsc{Union}$(x, y)$}{
  $u \gets \textsc{Find-Set}(x)$\;
  $v \gets \textsc{Find-Set}(y)$\;
  \lIf{$u = v$}{\Return \tcp*[h]{already in the same set; nothing to do}}
  $\textsc{Link}(u, v)$\;
}
\Fn{\textsc{Link}$(u, v)$}{
  \tcp{$u$, $v$ are distinct roots}
  \uIf{$u.\textit{rank} > v.\textit{rank}$}{
    $v.p \gets u$\;
  }
  \Else{
    $u.p \gets v$\;
    \If{$u.\textit{rank} = v.\textit{rank}$}{
      $v.\textit{rank} \gets v.\textit{rank} + 1$\;
    }
  }
}
\end{algorithm}

\begin{figure}[h]
\centering
\begin{tikzpicture}[
  every node/.style={draw, circle, minimum size=6.5mm, fill=blue!4,
    inner sep=0pt, font=\footnotesize},
  rt/.style={fill=green!12, very thick},
  >={Stealth[length=1.8mm]}, edge from parent/.style={draw,->},
  level distance=8mm, sibling distance=8mm]
  % Case 1: unequal ranks
  \begin{scope}
    \node[font=\footnotesize\scshape, text=gray!70!black, draw=none, fill=none] at (8mm, 15mm) {Case 1: Unequal Ranks};
    
    % Subtree triangles drawn first (behind nodes)
    \draw[fill=blue!6, draw=gray!60] (-5mm,-8mm) -- (5mm,-8mm) -- (0, -2mm) -- cycle;
    \draw[fill=blue!6, draw=gray!60] (11.5mm,-4.5mm) -- (20.5mm,-4.5mm) -- (16mm, 0mm) -- cycle;
    
    \node[rt] (u) at (0,0) {$u$};
    \node[rt] (v) at (16mm,2mm) {$v$};
    \draw[->, very thick] (v) -- (u);
    
    \node[draw=none,fill=none,font=\footnotesize] at (0,-11mm) {rank $2$};
    \node[draw=none,fill=none,font=\footnotesize] at (16mm,-7.5mm) {rank $1$};
    
    \node[draw=none,fill=none,font=\scriptsize,align=center] at (8mm,-18mm)
      {Lower rank $v$ hangs under $u$;\\root rank remains $2$.};
  \end{scope}
  % Case 2: equal ranks
  \begin{scope}[xshift=52mm]
    \node[font=\footnotesize\scshape, text=gray!70!black, draw=none, fill=none] at (7mm, 15mm) {Case 2: Equal Ranks};
    
    % Subtree triangles drawn first (behind nodes)
    \draw[fill=blue!6, draw=gray!60] (-4.5mm,-7.5mm) -- (4.5mm,-7.5mm) -- (0,-2mm) -- cycle;
    \draw[fill=blue!6, draw=gray!60] (9.5mm,-3.5mm) -- (18.5mm,-3.5mm) -- (14mm,2mm) -- cycle;
    
    \node[rt] (u2) at (0,0) {$u$};
    \node[rt] (v2) at (14mm,4mm) {$v$};
    \draw[->, very thick] (u2) -- (v2);
    
    \node[draw=none,fill=none,font=\footnotesize] at (0,-10.5mm) {was rank $1$};
    \node[draw=none,fill=none,font=\footnotesize] at (14mm,10mm) {rank $1{\to}2$};
    
    \node[draw=none,fill=none,font=\scriptsize,align=center] at (7mm,-18mm)
      {Ranks tied at $1$:\\one becomes root, its rank\\increments to $2$.};
  \end{scope}
\end{tikzpicture}
\caption{The two cases of \textsc{Link}($u, v$) in union by rank. When ranks differ, the smaller rank tree is attached under the larger with no rank change. When ranks match, either root may be chosen as parent, and its rank increases by exactly $1$.}
\label{ch19:fig-union-by-rank}
\end{figure}

The two cases of \textsc{Link}: a triangle is the rest of each root's tree.
Only the equal-rank case ever raises a rank, and only by one.

\subsection{Why rank controls height}
\label{ch19:rank-bound}

The payoff is a clean guarantee, and the argument for it explains exactly
\emph{why} rank works.

\begin{lemma}
\label{ch19:lem-rank-size}
With union by rank (and no path compression), any root of rank $r$ is the root
of a tree containing at least $2^r$ nodes.
\end{lemma}

\begin{proof}[Why it is true]
We argue by induction on the number of links performed. A fresh node has rank
$0$ and a tree of $1 = 2^0$ node, so the base case holds. Now consider a link.
If two roots of \emph{different} rank are joined, the surviving root's rank does
not change, and its tree only grew, so its size still beats $2^r$. The only
interesting case is linking two roots of \emph{equal} rank $r$. By the
inductive hypothesis each had at least $2^r$ nodes, so the merged tree has at
least $2^r + 2^r = 2^{r+1}$ nodes -- and its new rank is exactly $r+1$. The
bound holds.
\end{proof}

\begin{corollary}
\label{ch19:cor-rank-log}
Every rank is at most $\lfloor \log_2 n \rfloor$, where $n$ is the number of
elements. Hence every tree has height at most $\lfloor \log_2 n \rfloor$, and
$\textsc{Find-Set}$ runs in $O(\log n)$ time.
\end{corollary}

\begin{proof}[Why it is true]
A root of rank $r$ owns at least $2^r$ nodes, and no tree can own more than $n$
nodes, so $2^r \le n$, giving $r \le \log_2 n$. Since rank is an upper bound on
height, the height is at most $\log_2 n$ too.
\end{proof}

So union by rank alone already matches the $O(\log n)$ per-operation behavior of
weighted lists -- but now in a tree, which sets up the second trick. A trace
makes it concrete.

\begin{labbox}
\textbf{Trace union by rank.} Start with singletons $a,b,c,d,e,f,g,h$, each
rank $0$. Execute, in order:
\[
\textsc{Union}(a,b),\ \textsc{Union}(c,d),\ \textsc{Union}(a,c),\
\textsc{Union}(e,f),\ \textsc{Union}(a,e).
\]
\begin{enumerate}
  \item $\textsc{Union}(a,b)$: equal ranks $0$. Hang $a$ under $b$; now
        $b.\textit{rank}=1$. Tree: root $b$, child $a$.
  \item $\textsc{Union}(c,d)$: equal ranks $0$. Hang $c$ under $d$;
        $d.\textit{rank}=1$. Tree: root $d$, child $c$.
  \item $\textsc{Union}(a,c)$: roots are $b$ (rank $1$) and $d$ (rank $1$).
        Equal! Hang $b$ under $d$; $d.\textit{rank}=2$. Now $d$ roots all of
        $a,b,c,d$.
  \item $\textsc{Union}(e,f)$: equal ranks $0$. Hang $e$ under $f$;
        $f.\textit{rank}=1$.
  \item $\textsc{Union}(a,e)$: roots are $d$ (rank $2$) and $f$ (rank $1$).
        Different! Hang the lower, $f$, under $d$. No rank changes; $d$ stays
        rank $2$.
\end{enumerate}
Final tree rooted at $d$: children $c$, $b$, $f$; under $b$ is $a$; under $f$ is
$e$. The deepest node ($a$ or $e$) is $2$ steps from the root -- and indeed
$\textit{rank}(d) = 2$. Check that this $8$-node tree has height $2 \le
\log_2 8 = 3$, as the corollary promises.
\end{labbox}

\begin{center}
\begin{tikzpicture}[
  every node/.style={draw, circle, minimum size=6.5mm, fill=blue!4,
    inner sep=0pt, font=\small},
  rt/.style={fill=green!12, very thick},
  >={Stealth[length=2mm]}, edge from parent/.style={draw,->},
  level distance=11mm, level 1/.style={sibling distance=15mm}]
  \node[rt] (droot) {$d$}
    child {node {$b$} child {node {$a$}}}
    child {node {$c$}}
    child {node {$f$} child {node {$e$}}};
  \node[draw=none,fill=none,font=\footnotesize,right=4mm of droot] {rank $2$};
\end{tikzpicture}
\end{center}

The tree the trace builds. Its height is $2$, matching $\textit{rank}(d)=2$, and
comfortably under the $\log_2 8 = 3$ ceiling.

\section{Trick Two: Path Compression}
\label{ch19:compression}

Now the second idea, and it is a lovely one. Every time you call
$\textsc{Find-Set}(x)$, you walk a path from $x$ up to the root. That walk is
free information: along the way you visited a whole chain of nodes, and you
\emph{now know the root for every one of them}. So why throw that away? On the
way back, repoint every node you visited directly at the root. The next
$\textsc{Find-Set}$ on any of them is then a single hop. This is
\emph{path compression}.

\begin{figure}[h]
\centering
\begin{tikzpicture}[font=\small]
  % Panel (a): Before Find(c) - Deep lookup path
  \begin{scope}
    \node[font=\footnotesize\scshape, text=gray!70!black] at (0, 3.8) {(a) Before: Deep Search Path};
    \node[circle, draw=targetgreen, fill=targetgreen!20, very thick, minimum size=7.5mm] (r) at (0, 3.0) {$r$};
    \node[circle, draw=sourceblue, fill=sourceblue!10, thick, minimum size=7.5mm] (a) at (0, 2.0) {$a$};
    \node[circle, draw=sourceblue, fill=sourceblue!10, thick, minimum size=7.5mm] (b) at (0, 1.0) {$b$};
    \node[circle, draw=quotientorange, fill=quotientorange!25, very thick, minimum size=7.5mm] (c) at (0, 0.0) {$c$};

    \draw[-{Latex[length=2mm]}, thick] (a) -- (r);
    \draw[-{Latex[length=2mm]}, thick] (b) -- (a);
    \draw[-{Latex[length=2mm]}, thick] (c) -- (b);

    % Upward query traversal trace
    \draw[-{Latex[length=2mm]}, very thick, dashed, quotientorange] (c.east) to[bend right=25] (b.east);
    \draw[-{Latex[length=2mm]}, very thick, dashed, quotientorange] (b.east) to[bend right=25] (a.east);
    \draw[-{Latex[length=2mm]}, very thick, dashed, quotientorange] (a.east) to[bend right=25] (r.east);
    \node[font=\tiny, text=quotientorange!90!black, right=8pt of b] {walk $3$ hops};
  \end{scope}

  % Transition arrow
  \draw[-{Latex[length=3mm]}, very thick, quotientorange] (1.6, 1.5) -- 
    node[above, font=\tiny, text=quotientorange, align=center] {Unwind stack:\\$x.p \gets \textsc{Find}(x.p)$} (3.4, 1.5);

  % Panel (b): Repointing shortcuts
  \begin{scope}[xshift=5.2cm]
    \node[font=\footnotesize\scshape, text=gray!70!black] at (0, 3.8) {(b) Rewiring Shortcuts to Root};
    \node[circle, draw=targetgreen, fill=targetgreen!20, very thick, minimum size=7.5mm] (r2) at (0, 3.0) {$r$};
    \node[circle, draw=sourceblue, fill=sourceblue!10, thick, minimum size=7.5mm] (a2) at (-1.2, 1.2) {$a$};
    \node[circle, draw=sourceblue, fill=sourceblue!10, thick, minimum size=7.5mm] (b2) at (0, 1.2) {$b$};
    \node[circle, draw=quotientorange, fill=quotientorange!25, very thick, minimum size=7.5mm] (c2) at (1.2, 1.2) {$c$};

    \draw[-{Latex[length=2.2mm]}, very thick, targetgreen!80!black] (a2) -- (r2);
    \draw[-{Latex[length=2.2mm]}, very thick, targetgreen!80!black] (b2) -- (r2);
    \draw[-{Latex[length=2.2mm]}, very thick, targetgreen!80!black] (c2) -- (r2);

    \node[font=\tiny, text=targetgreen!90!black] at (0, 0.4) {All $3$ nodes repointed};
    \node[font=\tiny, text=gray!70!black] at (0, 0.05) {Direct parent is root $r$};
  \end{scope}

  % Summary annotation
  \node[draw=gray!30, fill=white, rounded corners=2pt, font=\scriptsize, align=left] at (9.6, 1.8) {
    \textbf{Path Compression Result}:\\
    $\bullet$ Path length shrinks from $3$ to $1$.\\
    $\bullet$ Subsequent finds on $a, b, c$\\
    \quad execute in exactly $1$ hop.\\
    $\bullet$ Work amortized across queries.
  };
\end{tikzpicture}
\caption{Path compression in \textsc{Find-Set}($c$). In (a), the search traverses the path $c \to b \to a \to r$. As the recursive calls return in (b), each visited node updates its parent pointer directly to root $r$. The tree is flattened, turning future lookups into instantaneous $O(1)$ operations.}
\label{ch19:fig-path-compression}
\end{figure}

A single $\textsc{Find}(c)$ walked $c \to b \to a \to r$; afterward $a$, $b$,
and $c$ \emph{all} point straight at $r$. The tree got dramatically flatter, and
it cost nothing extra -- we were walking that path anyway. The implementation is
a two-line recursion.

\begin{algorithm}[H]
\caption{\textsc{Find-Set} with path compression}
\SetKwProg{Fn}{function}{}{}
\Fn{\textsc{Find-Set}$(x)$}{
  \If{$x \neq x.p$}{
    $x.p \gets \textsc{Find-Set}(x.p)$\;
    \tcp{recurse to the root, then repoint $x$ straight at it}
  }
  \Return $x.p$\;
}
\end{algorithm}

Read that recursion carefully. The call $\textsc{Find-Set}(x.p)$ returns the
root and, as a side effect, compresses everything above $x$. Then the line
$x.p \gets (\text{that root})$ snaps $x$ itself onto the root. By the time the
outermost call returns, every node on the original path points directly at the
root.

\begin{trapbox}
Path compression changes node \emph{depths} but we deliberately do \emph{not}
update ranks when it happens. After compression a node's rank may be larger
than its actual height. That is fine and intentional: rank stays a valid
\emph{upper bound} on height, which is all our analysis ever needed, and
recomputing exact heights after every find would be slow and pointless. Rank is
a budget, not a measurement.
\end{trapbox}

\begin{walkbox}
Here is something delightful to chew on. Path compression makes
$\textsc{Find-Set}$ \emph{change the structure} even though it is, semantically,
a read-only query -- you only asked ``who is the representative?'' yet the tree
quietly rearranged. This is the same flavor of self-improvement you saw in
amortized analysis: an operation pays a little extra now to make all future
operations cheaper. Convince yourself that the answer it returns is never
wrong, no matter how much it rearranges -- the root of $x$'s tree is the same
before and after, because compression never moves a node to a \emph{different}
tree.
\end{walkbox}

\section{Putting It Together, and the Near-Constant Result}
\label{ch19:together}

Use both tricks -- union by rank \emph{and} path compression -- on the forest,
and the bound on a sequence of operations drops from $O(m \log n)$ to something
almost, but not quite, linear. Here is the celebrated result; afterward we will
work to make the strange function in it feel ordinary.

\begin{formalbox}
\textbf{Theorem (Tarjan).} On a forest with union by rank and path
compression, any sequence of $m$ disjoint-set operations, of which $n$ are
\textsc{Make-Set}, runs in total time
\[
O\big(m \, \alpha(n)\big),
\]
where $\alpha$ is the \emph{inverse Ackermann function}.
\end{formalbox}

The function $\alpha(n)$ grows so slowly that for every $n$ you will ever
encounter -- every $n$ that could be written down with atoms in the observable
universe -- $\alpha(n) \le 4$. So the average cost per operation is bounded by a
constant that, in any real program, never exceeds $4$ or $5$. The structure is
\emph{effectively} constant-time per operation, even though, famously, it is
provably \emph{not} constant in the strict asymptotic sense.

\subsection{Where does inverse Ackermann come from?}
\label{ch19:ackermann}

You do not need the full proof -- it is intricate, and CLRS devotes a careful
section to it. But you should know \emph{why} such a bizarre function shows up at
all. The whole effect comes from how violently fast a related function
\emph{grows}, and $\alpha$ is its sluggish inverse.

Start with a function that grows absurdly fast. Define a tower of operations
where each level is the iterate of the level below:

\begin{itemize}
  \item Repeated addition is multiplication: $a \cdot b = a + a + \cdots + a$
        ($b$ times).
  \item Repeated multiplication is exponentiation: $a^b = a \cdot a \cdots a$
        ($b$ times).
  \item Repeated exponentiation is \emph{tetration}: a tower
        $a^{a^{\cdot^{\cdot^{a}}}}$ of height $b$.
\end{itemize}

Keep climbing this ladder of operations -- each one iterating the previous --
and you get the \emph{Ackermann function} $A(k, j)$, which grows faster than any
primitive-recursive function. To get a feel for the violence, the iterated
power of two,
\[
2,\quad 2^2 = 4,\quad 2^{2^2} = 16,\quad 2^{2^{2^2}} = 65536,\quad
2^{2^{2^{2^2}}} = 2^{65536},
\]
reaches a number with about $20{,}000$ digits in just \emph{five} steps. The
inverse Ackermann function $\alpha(n)$ asks the reverse question:
\emph{how many steps up this tower do I need before I exceed $n$?} Because the
tower explodes so fast, the answer stays tiny forever.

\begin{center}
\begin{tikzpicture}[font=\small, >={Stealth[length=2mm]},
  tower/.style={draw, rounded corners, fill=structurepurple!12,
    minimum height=6mm, inner xsep=4pt}]
  \node[tower] (s1) {$2$};
  \node[tower, right=6mm of s1] (s2) {$4$};
  \node[tower, right=6mm of s2] (s3) {$16$};
  \node[tower, right=6mm of s3] (s4) {$65536$};
  \node[tower, right=6mm of s4] (s5) {$2^{65536}$};
  \draw[->] (s1)--(s2); \draw[->] (s2)--(s3);
  \draw[->] (s3)--(s4); \draw[->] (s4)--(s5);
  \node[font=\footnotesize, gray, below=2mm of s3]
    {each step raises $2$ to the previous value: explosive up, glacial inverse};
\end{tikzpicture}
\end{center}

Reading this strip \emph{left to right} is the tower exploding; reading it
\emph{right to left} -- ``how many boxes back from $n$ to reach $2$?'' -- is
$\alpha$. Five boxes already carry you past anything you could store, so
$\alpha(n)$ never climbs past a tiny handful.

\begin{intuitionbox}
Think of $\alpha(n)$ as the answer to: ``how many times must I apply this
ferociously-fast tower-building before I pass $n$?'' Since one application can
leap from $65536$ to $2^{65536}$, you essentially never need more than a
handful of applications. The single-variable inverse Ackermann
$\alpha(n) \le 4$ for all $n$ below roughly $2^{2^{2^{2^{16}}}}$ -- a number so
large it has no name and no physical meaning. For your laptop, your cluster,
and every dataset humanity will ever assemble, $\alpha(n)$ is a small constant.
\end{intuitionbox}

\subsection{A gentle sketch of why we win}
\label{ch19:sketch}

Here is the shape of the argument, without the full machinery. Union by rank
alone guarantees that ranks along any root path strictly increase, and that
ranks never exceed $\log_2 n$. Path compression then attacks those paths: every
time a node is found, it leaps closer to the root, so a node can only be ``far
from the root'' a limited number of times before it is essentially adjacent to
the root.

The accounting trick (this is amortized analysis, the previous chapter's tool)
groups nodes into rank \emph{blocks} and charges the cost of each step of a
find either to the operation itself or to a node whose rank is about to jump it
into a higher block. A node can only be charged this way a number of times
related to how many rank blocks exist -- and because ranks grow only to
$\log_2 n$ while the block boundaries are spaced by the tower function, the
number of blocks is about $\alpha(n)$. Multiply $m$ operations by $\alpha(n)$
charges each and you get the $O(m\,\alpha(n))$ bound.

\begin{focusbox}
You do not need to reproduce the proof. You need three things: \textbf{(1)} both
heuristics are needed for the $\alpha$ bound -- either one alone gives only
$O(\log n)$ per operation; \textbf{(2)} $\alpha$ is the inverse of a function
that grows faster than any tower, so it is bounded by a tiny constant for all
realistic $n$; \textbf{(3)} therefore union-find is, in every practical sense,
constant time per operation.
\end{focusbox}

\begin{historybox}
The disjoint-set problem was formalized in 1964 by Bernard A.~Galler and Michael J.~Fischer while designing compiler algorithms for FORTRAN's \lstinline|EQUIVALENCE| statement, which aliased variables in shared memory. Galler and Fischer introduced the disjoint-set forest representation and the union-by-rank and path-compression heuristics. In 1975, Robert E.~Tarjan published his celebrated analysis establishing the tight $\Theta(m\,\alpha(m,n))$ bound, which he later refined with Jan van Leeuwen (1984).

For over a decade, theorists assumed that path compression together with union by rank ran in strictly linear $O(m)$ time, suspecting that the exotic inverse-Ackermann factor $\alpha(n)$ was merely an artifact of an imperfect upper-bound proof. In a major theoretical breakthrough, Tarjan proved a matching lower bound: on any pointer machine where operations must follow pointers, \emph{no} disjoint-set algorithm can achieve $o(m\,\alpha(m,n))$ worst-case time. The inverse-Ackermann factor is not a mathematical defect; it reflects the inescapable cost of maintaining dynamic tree connectivity without cycles.

Tarjan's analysis proved that disjoint-set forests provide practically instantaneous dynamic connectivity checks. This enables near-linear running time for Kruskal's minimum spanning tree algorithm, connected-component labeling in computer vision, Hindley--Milner type unification in programming language compilers, and physical percolation modeling. Yet this carries an honest trade-off: monotonicity limitations. The standard structure only supports union operations; it cannot split sets, undo previous merges, or delete individual elements without reconstructing the forest from scratch. Designing fully dynamic connectivity algorithms that support both insertions and deletions of edges in polylogarithmic time remains one of the most vibrant areas of graph data structures research.
\end{historybox}

\subsection{What each variant costs}
\label{ch19:costs}

\begin{complexitybox}
For a sequence with $n$ \textsc{Make-Set}s and $m$ operations total:
\begin{center}
\begin{tabular}{lll}
\toprule
Representation & \textsc{Find-Set} (worst) & Sequence cost \\
\midrule
Linked list, arbitrary union & $O(1)$ & $O(n^2)$ in the bad case \\
Linked list, weighted union & $O(1)$ & $O(m + n \log n)$ \\
Forest, naive link & $O(n)$ & $O(mn)$ in the bad case \\
Forest, union by rank only & $O(\log n)$ & $O(m \log n)$ \\
Forest, path compression only & --- & $O(m \log n)$ (amortized) \\
Forest, rank \emph{and} compression & --- & $\boxed{O(m\,\alpha(n))}$ \\
\bottomrule
\end{tabular}
\end{center}
The last row is what you should actually implement. The two heuristics together
cost two extra fields per node (a parent pointer and a small rank) and a few
lines of code, and they buy you effectively constant time.
\end{complexitybox}

\section{A Complete Worked Example}
\label{ch19:fullexample}

Now run the finished structure -- forest, union by rank, path compression --
on a small connected-components problem so you see all the pieces move together.

Vertices: $1, 2, 3, 4, 5, 6$. Edges arrive in this order:
\[
(1,2),\ (3,4),\ (5,6),\ (1,3),\ (2,5).
\]
Every vertex starts as a singleton: parent itself, rank $0$.

\begin{labbox}
\textbf{Trace.} We write $p[\,]$ for the parent array and $r[\,]$ for ranks.
Initially $p[i]=i$ and $r[i]=0$ for all $i$.
\begin{enumerate}
  \item \textbf{Edge $(1,2)$.} $\textsc{Find}(1)=1$, $\textsc{Find}(2)=2$.
        Equal rank, link $1$ under $2$: $p[1]=2$, $r[2]=1$.
  \item \textbf{Edge $(3,4)$.} Roots $3,4$, equal rank: $p[3]=4$, $r[4]=1$.
  \item \textbf{Edge $(5,6)$.} Roots $5,6$, equal rank: $p[5]=6$, $r[6]=1$.
        Now three components: $\{1,2\}$, $\{3,4\}$, $\{5,6\}$.
  \item \textbf{Edge $(1,3)$.} $\textsc{Find}(1)$ walks $1\to 2$, returns $2$
        (rank $1$). $\textsc{Find}(3)$ walks $3\to 4$, returns $4$ (rank $1$).
        Equal rank, link $2$ under $4$: $p[2]=4$, $r[4]=2$. Component
        $\{1,2,3,4\}$ rooted at $4$.
  \item \textbf{Edge $(2,5)$.} $\textsc{Find}(2)$ walks $2\to 4$, returns $4$
        (rank $2$). $\textsc{Find}(5)$ walks $5\to 6$, returns $6$ (rank $1$).
        Different rank: hang the smaller, $6$, under $4$: $p[6]=4$. No rank
        change. Everything is now one component, rooted at $4$.
\end{enumerate}
\textbf{Now watch path compression flatten the tree.} Call $\textsc{Find}(5)$.
The path is $5 \to 6 \to 4$. Compression repoints: $p[5]=4$ and $p[6]=4$ (it
already was). Next call $\textsc{Find}(1)$: path $1 \to 2 \to 4$; afterwards
$p[1]=4$ and $p[2]=4$. After a few finds, \emph{every} vertex points directly
at $4$, and all subsequent finds are a single hop. The tree has been ironed
flat.
\end{labbox}

\begin{center}
\begin{tikzpicture}[
  every node/.style={draw, circle, minimum size=6.5mm, fill=blue!4,
    inner sep=0pt, font=\small},
  >={Stealth[length=2mm]},
  level 1/.style={sibling distance=20mm},
  level distance=12mm,
  edge from parent/.style={draw,->}]
  \node[fill=green!12, very thick] {$4$}
    child {node {$1$}}
    child {node {$2$}}
    child {node {$3$}}
    child {node {$5$}}
    child {node {$6$}};
  \node[draw=none, fill=none, below=18mm, font=\small]
    {after compression: a flat star, every \textsc{Find} is one hop};
\end{tikzpicture}
\end{center}

\section{Implementation Notes}
\label{ch19:impl}

In real code you almost never use explicit node objects with pointers. You use
two integer arrays indexed $0 \dots n-1$: a \texttt{parent} array and a
\texttt{rank} array. $\textsc{Make-Set}(i)$ sets \lstinline|parent[i] = i| and
\lstinline|rank[i] = 0|. This is cache-friendly and executes in direct $O(1)$ time.

A common and excellent variant of $\textsc{Find}$ avoids recursion (and its
stack depth) using \emph{path halving}: as you walk up, make every node point to
its grandparent. It compresses almost as well as full compression and runs in a
single tight loop.

\begin{algorithm}[H]
\caption{\textsc{Find-Set} with path halving (iterative)}
\SetKwProg{Fn}{function}{}{}
\Fn{\textsc{Find-Set}$(x)$}{
  \While{$x \neq \texttt{parent}[x]$}{
    $\texttt{parent}[x] \gets \texttt{parent}[\texttt{parent}[x]]$\;
    \tcp{point $x$ at its grandparent, then step up}
    $x \gets \texttt{parent}[x]$\;
  }
  \Return $x$\;
}
\end{algorithm}

\begin{trapbox}
A subtle but classic bug: comparing the \emph{representatives} but unioning the
\emph{original elements} correctly. In $\textsc{Union}(x,y)$ you must call
$\textsc{Find-Set}$ on both \emph{first}, link the two roots, and never link an
element to itself. Linking a root to itself ($u.p \gets u$ when $u = v$) is a
silent no-op that wastes a rank increment in some buggy variants; always check
that the two roots are distinct before incrementing any rank.
\end{trapbox}

Many implementations replace rank with \emph{size} (the number of nodes in the
tree) and link the smaller tree under the larger -- ``union by size.'' It
achieves the same $O(\alpha(n))$ bound, and size is sometimes useful on its own
(``how big is this component?''). Rank and size are two flavors of the same
``small under big'' idea you have now seen three times: weighted lists, union by
rank, union by size. Whenever you merge structures, absorbing the small into the
large is almost always the move.

\section{Summary}
\label{ch19:summary}

You started with a problem so plain it barely seemed to need an algorithm:
group objects, merge groups, ask if two objects share a group. You built it four
ways, each better than the last.

\begin{itemize}
  \item \textbf{Linked lists} give instant $\textsc{Find}$ but slow
        $\textsc{Union}$; the weighted-union heuristic (small list onto big)
        tames it to $O(m + n \log n)$.
  \item \textbf{Forests} name each set by a tree root. Naive linking can build
        a stringy path of height $n$.
  \item \textbf{Union by rank} attaches the shorter tree under the taller,
        keeping height $O(\log n)$.
  \item \textbf{Path compression} flattens the find path on every query, so
        the structure improves itself as it is used.
\end{itemize}

Together, rank and compression give $O(m\,\alpha(n))$ -- effectively constant
time per operation, where $\alpha$, the inverse Ackermann function, is bounded
by a tiny constant for every input you will ever meet. Two integer arrays and a
dozen lines of code buy you one of the fastest data structures in existence.

\begin{aestheticbox}
Notice the design pattern behind both heuristics: each operation invests a small
amount of work up front to make \emph{future} operations cheaper. Weighted union
bounds how often any element can change leaders, and path compression flattens
pointer chains lazily during find queries, postponing structural cleanup until
someone actually visits those nodes.
\end{aestheticbox}

In \textbf{Chapter~20 (Fibonacci Heaps)}, we apply this philosophy of lazy postponement
to priority queues. Instead of reorganizing heaps on every key update, a Fibonacci heap
does almost nothing during insertions and key decreases, letting roots accumulate in a loose collection.
Only when \textsc{Extract-Min} demands the smallest item does the structure perform a sweeping,
consolidating cleanup. Using the potential method of amortized analysis, we will prove that this
procrastination pays off: insertions, merges, and key decreases all run in $O(1)$
amortized time, accelerating graph algorithms like Dijkstra and Prim.

\section{Exercises}
\label{ch19:exercises}

\begin{exercise}
\textbf{(Warmup: trace it.)} Start with singletons $1$ through $8$. Using the
forest with union by rank (no path compression), perform $\textsc{Union}(1,2)$,
$\textsc{Union}(3,4)$, $\textsc{Union}(5,6)$, $\textsc{Union}(7,8)$,
$\textsc{Union}(1,3)$, $\textsc{Union}(5,7)$, $\textsc{Union}(1,5)$. Draw the
final tree and give the rank of the root. What is the height of the tree?
\end{exercise}

\begin{exercise}
\textbf{(Same set?)} In one or two sentences, explain why the test
$\textsc{Find-Set}(x) = \textsc{Find-Set}(y)$ correctly decides whether $x$ and
$y$ are in the same set, and why no separate operation is needed. What property
of representatives makes this work?
\end{exercise}

\begin{exercise}
\textbf{(The weighted-union doubling argument.)} In the linked-list
representation with weighted union, prove carefully that a single fixed object's
back-pointer is updated at most $\lfloor \log_2 n \rfloor$ times over any
sequence of operations. Where exactly does the ``set at least doubles'' step
come from?
\end{exercise}

\begin{exercise}
\textbf{(Rank versus height.)} Give a concrete sequence of operations (using
both union by rank and path compression) after which some node's rank is
strictly greater than the actual height of its tree. Explain why this does not
break the analysis.
\end{exercise}

\begin{exercise}
\textbf{(Path compression alone.)} Suppose you use path compression but
\emph{not} union by rank -- $\textsc{Link}$ always hangs the first root under
the second. Argue informally that the amortized cost per operation is still
$O(\log n)$, not the bad $O(n)$ of the naive forest. (Hint: compression keeps
ironing the trees flat even without rank.)
\end{exercise}

\begin{exercise}
\textbf{(Union by size.)} Rewrite $\textsc{Make-Set}$, $\textsc{Link}$, and the
rank logic to use ``union by size'' instead: store the number of nodes in each
tree at its root, and on a link hang the smaller-size tree under the larger.
Prove the analogue of Lemma~\ref{ch19:lem-rank-size}: a tree of height $h$
contains at least $2^h$ nodes, so heights stay $O(\log n)$.
\end{exercise}

\begin{exercise}
\textbf{(Counting components.)} Augment the structure with a global counter so
that after any sequence of operations you can report, in $O(1)$ time, how many
disjoint sets currently exist. Specify exactly where the counter is initialized,
incremented, and decremented, and prove it stays correct.
\end{exercise}

\begin{exercise}
\textbf{(Kruskal preview.)} You are given a graph whose edges arrive one at a
time, and you must print ``CYCLE'' the first time an edge connects two vertices
already in the same component, and ``ADD'' otherwise. Write pseudocode using
union-find, and argue its total running time over $m$ edges and $n$ vertices is
$O((m+n)\,\alpha(n))$. You have just written the core of Kruskal's algorithm.
\end{exercise}

\begin{exercise}
\textbf{(Scale of $\alpha$, really?)} The two-argument Ackermann-style tower
satisfies (roughly) $A(2) = 4$, $A(3) = 16$, $A(4) = 65536$, and
$A(5) = 2^{65536}$. Using the intuition that $\alpha(n)$ is ``the least $k$ with
$A(k) \ge n$,'' explain why $\alpha(n) \le 4$ for every $n$ that could be stored
in the memory of any physically possible computer. Why does this make the
``not constant time'' result a mathematical curiosity rather than a practical
worry?
\end{exercise}

\begin{exercise}
\textbf{(Challenge: offline minimum / interval merging.)} You process a stream
of integers from a fixed range $1 \dots n$. Some are ``insert $k$'' and some are
``extract the current minimum.'' Design a way to use union-find so that each
extract finds the next still-present value quickly, by unioning each removed
value with its successor. Explain why each value is removed at most once and the
whole process is $O((n + q)\,\alpha(n))$ for $q$ operations. (This is the
classic ``union-find as a pointer to the next free slot'' trick.)
\end{exercise}
